\section{Symmetry: using a group without enumerating it}
\label{sec:symmetry}

A permutation group on $n$ points can have $n!$ elements. A proof planner that
wants to exploit symmetry --- to merge equivalent cases, to count orbits, to
reduce a cone search --- cannot enumerate them, and one contributing effort
supplies the certificate language that makes the group usable without doing so.
This is the collection's only group theory.

\subsection{The certificate, and what makes non-membership checkable}

The object is a \emph{stabiliser chain}. Fix the base $0,1,\ldots,n-1$ and let
$K_i$ be the subgroup fixing the first $i$ points. The certificate supplies,
per level, the orbit of the $i$th point under $K_i$ as a \emph{tree} --- each
non-root entry recording its point, an earlier parent, and the signed generator
reaching it --- together with a claimed group order and a word DAG deriving
every strong generator from the \emph{original} supplied generators.

The claimed order is checked by the product $\prod_i|\Omega_i|$, and the effort
is blunt that this check alone is worthless: ``a malformed collection of tables
also has a product.'' What makes the product a group order is the Schreier
condition, and the acceptance requirement is that \emph{every} Schreier residue
$t_{s(x)}^{-1}s\,t_x$ sifts to the identity --- every generator, every orbit
point, no sampling and no producer-chosen subset.

\begin{theorem}[Complete chain]\label{thm:chain}
An accepted certificate yields a bijection
$T_0\times\cdots\times T_{n-1}\to G$ by multiplication. Consequently the order
product is $|G|$, and sifting decides membership.
\end{theorem}
\begin{proof}
The strong list generates $G$, so $K_0=G$, and the Schreier condition
identifies each $K_{i+1}$ with the stabiliser of $i$ in $K_i$. The bottom group
is trivial because a permutation fixing every point is the identity, so
repeated coset decomposition gives surjectivity. For uniqueness, compare two
products: their values at $0$ select the same point of $\Omega_0$ and hence the
same stored transversal; cancel and continue down the base. If a member of $G$
failed sifting, its unique coset digit would be absent from an exact orbit,
which contradicts the orbit being exactly $K_i\cdot i$.
\end{proof}

That last sentence is the answer to the question this section exists to
answer. A failed sift is a \emph{negative} membership witness --- and it is one
only because the chain's completeness was established first, by a finite,
fully enumerated check. This is not a bounded search reporting exhaustion. It
is a decision, and the finiteness is the number of Schreier edges rather than a
budget.

The regression fixture that isolates the gate is the right one: take
$\Sym_3$, supply all three orbit points at level $0$ but only the roots below,
and the order product is correct, every listed orbit is closed under its listed
generators, and both original generators are present. The certificate is still
rejected, because a residue fails to sift. Every superficial check passes and
the logical gate catches it.

\subsection{Canonical images, and a minimum that is not the sorted tuple}

For symmetry breaking the useful operation is the canonical image: the
lexicographically least $g\cdot c$ over $g\in G$. It is certified by a coset
cover --- a tree whose splits enumerate every transversal at a level and whose
cuts carry a proved lower bound --- plus one transporter attaining the claimed
minimum. Coverage proves no coset was skipped; the transporter proves the bound
is reached.

The worked case that repays attention is the alternating group. For
$\Alt_n$ the canonical image of a tuple with all-distinct entries is
\emph{not} the sorted tuple when the sorting permutation is odd: it is the
sorted tuple with its last two entries exchanged. With a repeated entry the
sorted tuple is reachable after all, because a transposition of two
equal-valued positions is an odd element of the stabiliser. Both branches are
in the corpus, checked against full enumeration at small degree.

\begin{remark}[What a symmetry tactic may and may not do]
The soundness conditions are exactly two equivalences,
\[
 \exists x\,P(x)\iff\exists x\,[P(x)\wedge r(x)=x],\qquad
 \forall x\,P(x)\iff\forall x\,[r(x)=x\Rightarrow P(x)],
\]
for a representative map $r$ and a $G$-invariant $P$. What they do \emph{not}
license is proving $x=r(x)$, rewriting occurrences of $x$ by $r(x)$ inside an
arbitrary expression, or conjoining a canonical-form constraint to the original
goal as though it held pointwise. Symmetry breaking preserves existence, or a
properly transported universal; it does not preserve the solution set.
\end{remark}

A second scope condition is sharper than it looks. A \emph{proper} subgroup of
the true symmetry group is perfectly sound for case merging --- it merges
fewer cases. It is not sound for orbit counting, where the answer depends on
the whole group. Both computations can be internally correct while answering
different questions, and only the source authorization distinguishes them.

\subsection{Counting, and a reduction for an existing worker}

Burnside's lemma turns the group into a polynomial: with $a_G(\lambda)$
elements of cycle type $\lambda$, the number of colourings up to symmetry is
$|G|^{-1}\sum_\lambda a_G(\lambda)q^{|\lambda|}$, universally in $q$. The
inventory is the certificate, and the checker recomputes it a different way ---
enumerating unique transversal normal forms from Theorem~\ref{thm:chain} and
recomputing cycles --- rather than re-running the producer's Cayley-graph walk.

Recognising $\Sym_n$ or $\Alt_n$ exactly, by order and parity gates rather than
by a probabilistic test, then replaces the inventory entirely by a closed
formula. That is what lets a degree-40 symmetric group be used: $40!$ is
certified in under twelve kilobytes of chain data, and the three-colour orbit
count follows from a binomial identity rather than from any enumeration.

The lane that matters most for the rest of this document is the Reynolds
projection. Averaging a polynomial over $G$ is a projection onto the invariants,
and on a monomial orbit it is computed without touching $|G|$ at all, because
each source monomial reaches every destination in its orbit the same number of
times. A cone search over a $G$-stable dictionary can therefore be replaced by
a search over orbit sums, and the effort's example collapses 780 coefficient
variables to one. The accepted reduced solution is then expanded back and
validated by the \emph{existing} cone checker of Section~\ref{sec:nonlinear} ---
an accelerator that does not move the trust boundary.

\begin{remark}[Averaging does not preserve the goal]
The projection is a tool for the search, not a substitute for the target. With
$G$ the swap and $p=x_0-x_1$, the average is zero, so $\mathcal{R}_G(p)\ge0$ holds
while $p\ge0$ does not. A tactic must not replace a goal by its average.
\end{remark}

\subsection{Scope and a formalisation hazard}

The fragment is concrete finite actions on $\{0,\ldots,n-1\}$ with permutations
given as arrays. Finitely presented groups, actions on infinite types, and
inferring a group from a program's behaviour are outside it. Representing an
abstract group by permutations needs its own faithfulness theorem before the
image's order may be called the group's order.

One Mathlib detail is worth carrying forward because it would silently
invalidate every counting certificate. The natural translation of ``number of
cycles'' via the length of \code{cycleType} undercounts: \code{cycleType}
omits fixed points, so the identity would be recorded as having zero cycles
rather than $n$. The correct exponent is
$\#\operatorname{cycleType}(g)+n-|\operatorname{supp}(g)|$. This is a small
representation mismatch of exactly the kind Section~\ref{sec:trust} exists to
catch, and it was found by reading the library rather than by running anything.
